\section{Source samples and averaging at nonempty occurrences}
\label{sec:peps_nonempty_source_samples}

\begin{definition}[Source samples at nonempty gate occurrences]
    \label{def:peps_nonempty_source_samples}
    \lean{TNLean.PEPS.PairEffect.SourceCircuit.sourceLocationsNonemptyEquiv,
      TNLean.PEPS.PairEffect.SourceCircuit.sourceSampleNonemptyEquiv}
    \leanok
    An actual source slot is labelled by its original gate occurrence
    and its slot within that gate. Its two recorded endpoints witness
    that the participating set of this occurrence is nonempty. The
    source slot therefore corresponds to its unordered endpoint pair
    at this nonempty original occurrence.
\end{definition}

\begin{theorem}[Preservation of the original occurrence and endpoints]
    \label{thm:peps_nonempty_source_sample_endpoints}
    \lean{TNLean.PEPS.PairEffect.SourceCircuit.sourceSampleNonemptyEquiv_gate,
      TNLean.PEPS.PairEffect.SourceCircuit.sourceSampleNonemptyEquiv_parties}
    \leanok
    \uses{def:peps_nonempty_source_samples}
    This correspondence is a bijection. It preserves the original
    gate occurrence and precisely the set of its two recorded endpoints.
    It introduces no source slot at an empty-participant occurrence.
\end{theorem}

\begin{proof}\leanok
    The slot-to-pair correspondence at each original gate is a bijection.
    Its inverse retains the gate and recovers the actual slot. Restricting
    the gate index to nonempty participating sets loses no source slot,
    since its left endpoint is a participating party.
\end{proof}

\begin{theorem}[Counting the actual source slots]
    \label{thm:peps_nonempty_source_slot_count}
    \lean{TNLean.PEPS.PairEffect.SourceCircuit.slotCount_eq_zero_of_not_nonempty,
      TNLean.PEPS.PairEffect.SourceCircuit.length_partialSlots_all_eq_sum_slotCount,
      TNLean.PEPS.PairEffect.SourceCircuit.length_partialSlots_all_eq_sum_nonempty_slotCount}
    \leanok
    \uses{def:peps_nonempty_source_samples}
    Let $n_g$ be the source-slot count of original occurrence $g$,
    and let $N$ be the total all-party source-slot count. An
    empty-participant occurrence has $n_g=0$, and
    $N=\sum_g n_g=\sum_{g:\,C_g\ne\varnothing}n_g$.
\end{theorem}

\begin{proof}\leanok
    The all-party source slots are in bijection with the pairs consisting
    of an original occurrence and one of its slots. Count this disjoint
    union, first over all occurrences and then over the nonempty ones.
\end{proof}

\begin{theorem}[Product of the averaging weights]
    \label{thm:peps_source_averaging_product}
    \lean{TNLean.PEPS.PairEffect.SourceCircuit.prod_sourceSlotWeight_eq}
    \leanok
    \uses{thm:peps_nonempty_source_slot_count}
    For every $z\in\mathbb C$, including $z=0$,
    $\prod_{g:\,C_g\ne\varnothing}z^{n_g}=z^N$.
    Thus, for a positive sample count $k$, the weight $k^{-n_g}$ at each nonempty original occurrence
    gives exactly the total source averaging factor $k^{-N}$.
\end{theorem}

\begin{proof}\leanok
    Add the exponents in the finite product and apply the slot-count identity.
\end{proof}

These correspondences and weights are those in
\cite[proof of Theorem~5.2]{OpenAI2026PolynomialPEPS},
\texttt{04-compression.tex}, lines 342--355 and 565--588.
